% ============================================================
% Cauchy Sequences — Canonical Proof Patterns
% ============================================================

\section{Canonical Cauchy Proof Patterns}

Cauchy sequences are one of the central organizing ideas of analysis.
A remarkable fact --- one that practicing analysts quickly internalize ---
is that the majority of Cauchy arguments follow a small number of
recognizable patterns.

The goal of this section is not merely to present results, but to train
the reader to recognize these recurring proof structures.

Almost every Cauchy proof in metric spaces ultimately reduces to two ideas:

\begin{center}
\textbf{Triangle inequality + tail control}
\end{center}

Symbolically, we repeatedly exploit estimates of the form

\[
d(x_m,x_n)
\le d(x_m,a)+d(a,x_n)
\]

for a carefully chosen \emph{anchor point} \(a\).

\bigskip

The following ten patterns occur again and again throughout analysis.

% ============================================================
\subsection*{Pattern 1 — Convergent $\Rightarrow$ Cauchy}
% ============================================================

\begin{remark*}[Techniques]
tech:triangle-inequality, tech:epsilon-splitting, tech:choose-N
\end{remark*}

\textbf{Claim.}
If \(x_n \to x\) in a metric space, then \((x_n)\) is Cauchy.

\begin{proof}

Let \(\varepsilon>0\) be given.

Because \(x_n\to x\), there exists \(N\) such that

\[
d(x_n,x) < \frac{\varepsilon}{2}
\qquad n\ge N
\]

Now take \(m,n\ge N\). By the triangle inequality,

\[
d(x_m,x_n)
\le d(x_m,x) + d(x,x_n).
\]

Each term is less than \(\varepsilon/2\), so

\[
d(x_m,x_n) < \varepsilon.
\]

Thus the sequence is Cauchy.
\end{proof}

\bigskip

Professor's advice:

\begin{quote}
Whenever a limit is known, you should immediately think of the triangle
inequality centered at the limit.
\end{quote}

\bigskip

% ============================================================
\subsection*{Pattern 2 — Cauchy $\Rightarrow$ Bounded}
% ============================================================

\begin{remark*}[Techniques]
tech:choose-N, tech:tail-argument, tech:ball-containment
\end{remark*}

\textbf{Idea.}

A Cauchy sequence eventually lives inside a small ball.

\textbf{Example.}

Let \((x_n)\) be Cauchy.

Choose \(N\) such that

\[
d(x_m,x_n) < 1
\quad
m,n\ge N.
\]

Fix \(n=N\). Then

\[
d(x_m,x_N)<1
\quad m\ge N.
\]

Hence the tail lies in the ball

\[
B(x_N,1).
\]

The initial segment

\[
\{x_1,\dots,x_{N-1}\}
\]

is finite, so the sequence is bounded.

\bigskip

\begin{quote}
A useful mental picture is that a Cauchy sequence eventually gets trapped
inside smaller and smaller balls.
\end{quote}

% ============================================================
\subsection*{Pattern 3 — Subsequences of Cauchy Sequences}
% ============================================================

\begin{remark*}[Techniques]
tech:subsequence-extraction, tech:tail-argument
\end{remark*}

If a sequence is Cauchy then any subsequence is also Cauchy.

Indeed, if

\[
d(x_m,x_n) < \varepsilon
\quad
m,n\ge N
\]

then any subsequence indices \(n_k,n_\ell\) eventually satisfy

\[
n_k,n_\ell \ge N.
\]

Therefore

\[
d(x_{n_k},x_{n_\ell}) < \varepsilon.
\]

\bigskip

This simple observation becomes extremely important in compactness and
Bolzano--Weierstrass arguments.

% ============================================================
\subsection*{Pattern 4 — Cauchy + Convergent Subsequence}
% ============================================================

\begin{remark*}[Techniques]
tech:triangle-inequality, tech:subsequence-extraction
\end{remark*}

Suppose

\[
x_{n_k} \to L
\]

and \((x_n)\) is Cauchy.

We estimate

\[
d(x_n,L)
\le
d(x_n,x_{n_k}) + d(x_{n_k},L).
\]

Choose \(k\) large so both terms are small.

Thus \(x_n \to L\).

\bigskip

\begin{quote}
This argument is the core mechanism behind many completeness proofs.
\end{quote}

% ============================================================
\subsection*{Pattern 5 — Algebra of Cauchy Sequences}
% ============================================================

\begin{remark*}[Techniques]
tech:triangle-inequality
\end{remark*}

Example: sum of Cauchy sequences.

\[
|a_m+b_m-(a_n+b_n)|
\le
|a_m-a_n|+|b_m-b_n|
\]

Each term becomes small independently.

Hence sums of Cauchy sequences are Cauchy.

% ============================================================
\subsection*{Pattern 6 — Metric Comparison}
% ============================================================

\begin{remark*}[Techniques]
tech:metric-verification
\end{remark*}

Suppose

\[
d_1(x,y) \le C\, d_2(x,y).
\]

Then any \(d_2\)-Cauchy sequence is also \(d_1\)-Cauchy.

Such estimates appear frequently when comparing metrics.

% ============================================================
\subsection*{Pattern 7 — Telescoping Argument}
% ============================================================

\begin{remark*}[Techniques]
tech:tail-argument
\end{remark*}

If

\[
x_n = \sum_{k=1}^{n} a_k
\]

then

\[
x_m-x_n
=
\sum_{k=n+1}^{m} a_k.
\]

Bounding the tail sum proves the sequence is Cauchy.

% ============================================================
\subsection*{Pattern 8 — Completeness Mechanism}
% ============================================================

\begin{remark*}[Techniques]
tech:subsequence-extraction, tech:triangle-inequality
\end{remark*}

Typical completeness proof structure:

\begin{enumerate}
\item Start with a Cauchy sequence
\item Prove it is bounded
\item Apply Bolzano--Weierstrass
\item Extract convergent subsequence
\item Apply Pattern 4
\end{enumerate}

% ============================================================
\subsection*{Pattern 9 — Non-completeness Counterexample}
% ============================================================

\begin{remark*}[Techniques]
tech:counterexample
\end{remark*}

Example in \(\mathbb{Q}\):

Let \(x_n\) be decimal truncations of \(\sqrt{2}\).

Then \(x_n\) is Cauchy in \(\mathbb{Q}\) but does not converge in
\(\mathbb{Q}\).

This demonstrates that \(\mathbb{Q}\) is not complete.

% ============================================================
\subsection*{Pattern 10 — Nested Interval Construction}
% ============================================================

\begin{remark*}[Techniques]
tech:diameter-bound
\end{remark*}

Suppose

\[
I_1 \supset I_2 \supset I_3 \supset \cdots
\]

and

\[
\text{diam}(I_n) \to 0.
\]

Choose \(x_n \in I_n\).

Then

\[
|x_m-x_n| \le \text{diam}(I_n)
\]

for \(m>n\), so the sequence is Cauchy.

\bigskip

% ============================================================
\section{Seven Canonical Triangle-Inequality Manipulations}
% ============================================================

The triangle inequality is the single most important inequality in
analysis.

Experienced analysts use a small repertoire of standard manipulations.

\bigskip

\subsection*{1. Insert an Intermediate Point}

\[
d(x,y)
\le
d(x,a)+d(a,y)
\]

Used to compare two quantities via an anchor point.

\bigskip

\subsection*{2. $\varepsilon$-Splitting}

\[
\varepsilon
=
\frac{\varepsilon}{2}
+
\frac{\varepsilon}{2}
\]

Used when two independent bounds must be controlled.

\bigskip

\subsection*{3. Three-Term Split}

\[
d(x,z)
\le
d(x,y)+d(y,w)+d(w,z)
\]

Useful in subsequence arguments.

\bigskip

\subsection*{4. Reverse Triangle Inequality}

\[
|\,|x|-|y|\,|
\le
|x-y|
\]

Often used to control absolute values.

\bigskip

\subsection*{5. Bounding a Difference by a Sum}

\[
|a-b|
\le
|a|+|b|
\]

Used constantly in algebraic estimates.

\bigskip

\subsection*{6. Telescoping Bound}

\[
|x_m-x_n|
\le
\sum_{k=n+1}^{m} |x_k-x_{k-1}|
\]

Key tool in series arguments.

\bigskip

\subsection*{7. Metric Ball Inclusion}

If

\[
d(x,a)<r/2
\quad\text{and}\quad
d(a,y)<r/2
\]

then

\[
d(x,y)<r.
\]

This observation underlies many compactness arguments.

\bigskip

\begin{center}
\textbf{Final Advice}
\end{center}

Most analysis proofs can be mentally reduced to the following question:

\begin{quote}
Where should I place the triangle inequality?
\end{quote}

Once the correct anchor point is chosen, the remainder of the argument
typically follows automatically.